\chapter*{Conclusion Générale}
\addcontentsline{toc}{chapter}{CONCLUSION GÉNÉRALE}
\adjustmtc
\thispagestyle{MyStyle}

Ce stage de fin d'études m'a permis de conduire un projet logiciel complet, depuis l'analyse du besoin jusqu'à la structuration d'une solution multi-interface prête à être démontrée et consolidée. Le système développé pour \textbf{FRS} répond à un besoin concret de centralisation de la relation client et propose une vision unifiée des activités commerciales, du support et des interactions avec les clients.\par
\vspace{0.3cm}

\textbf{Bilan du travail réalisé.} Le projet a été mené de manière progressive à travers quatre sprints complémentaires :
\begin{itemize}
    \item le premier sprint a permis d'étudier l'existant, d'identifier les acteurs et de formaliser les besoins ;
    \item le deuxième sprint a servi à concevoir les modèles de données, les diagrammes UML, l'architecture technique et l'API REST ;
    \item le troisième sprint a été consacré à l'implémentation du backend, du frontend, de l'application mobile, du portail client et du module d'assistance documentaire ;
    \item le quatrième sprint a porté sur la validation fonctionnelle, la revue des permissions, la correction des anomalies de stabilisation et la préparation du déploiement local via Docker Compose.
\end{itemize}

\vspace{0.3cm}
\textbf{Résultats obtenus.} La solution développée fournit aujourd'hui :
\begin{itemize}
    \item un cycle commercial structuré autour des leads, comptes, contacts, deals et devis ;
    \item un module de support avec tickets, statuts, priorités et messagerie publique ou interne ;
    \item un portail client web sécurisé pour le suivi des tickets et des devis ;
    \item une application mobile client permettant l'accès nomade aux fonctionnalités essentielles ;
    \item des tableaux de bord de pilotage pour les ventes et le support ;
    \item un journal d'audit pour assurer la traçabilité des opérations sensibles ;
    \item une base documentaire interrogeable via un assistant conversationnel fondé sur la technique RAG ;
    \item une infrastructure locale conteneurisée facilitant l'exécution du projet et de ses campagnes de tests.
\end{itemize}

\vspace{0.3cm}
\textbf{Apports techniques et professionnels.} Ce projet m'a permis de consolider plusieurs compétences complémentaires : conception d'API avec FastAPI, modélisation de données relationnelles, développement d'interfaces web avec React, développement mobile avec React Native, gestion de l'authentification JWT, séparation des responsabilités par rôles, intégration d'un moteur de recherche documentaire assisté par IA, et organisation d'un projet logiciel selon une logique itérative. Il m'a également appris à relier des besoins métier concrets à des choix d'architecture réalistes et maintenables.\par
\vspace{0.3cm}

\textbf{Perspectives.} Plusieurs axes d'amélioration peuvent prolonger ce travail :
\begin{itemize}
    \item enrichir le portail client par des notifications temps réel ;
    \item renforcer la stratégie de migration et d'administration de la base de données ;
    \item étendre la couverture de tests automatisés sur l'ensemble du cycle d'exécution ;
    \item connecter l'assistant à des sources de données métiers plus dynamiques selon des règles de sécurité strictes ;
    \item préparer un déploiement cloud industrialisé avec supervision et gestion des secrets.
\end{itemize}

En définitive, ce stage a constitué une expérience formatrice et professionnalisante. Il m'a donné l'occasion de participer à la construction d'une solution cohérente, utile et techniquement riche, tout en renforçant ma capacité à concevoir des systèmes logiciels réalistes, documentés et évolutifs.\par
